\section{Introduction}
\label{sec:intro}

An agent built on a large language model works by calling tools and reading what they return. The tools are
usually written for programs, not for models: a keyword-research API returns monthly volumes, click-through data,
backlink counts and a dozen normalisations of each; a payment API returns every field of a charge, null or not; a
code-hosting API embeds a full user object, with fifteen navigation URLs, in every issue. All of it enters the
context, is paid for at every subsequent turn, and competes for the model's attention with what the task needs.

Practitioners handle this by hand: a wrapper function per tool that keeps the fields the task uses. This works,
but it requires knowing which fields those are, for each tool and each task, and keeping the wrappers in step
with the APIs. Generic compressors avoid the per-tool work, but they mostly remove structural redundancy; those that learn
per tool learn from what the model asks back, not from what the task's answers use.

This paper asks whether the hand-written selection can be \emph{learned}. The agent's own execution traces show
what it did with each tool output: which values reappear in the arguments of later calls and in the final answer.
A field whose values never reappear, across runs of the same task, is a candidate for removal. We call the
resulting filter a \emph{learned projection}, implement it in an open-source library (\texttt{tracetrim}), and
test it on five agents: three that use real APIs (an editorial planner and a content-gap analyst on a search-data
API, a literature scanner on a scholarly catalogue) and two that use simulated servers whose payloads follow the
public OpenAPI descriptions of a payment API and a code-hosting API.

Since version 0.4, the library protects identifiers, recognises transformed values, writes compact JSON,
offers \texttt{expand} for any hidden field and can keep learning while the agent works. Measured with these
defaults, in live runs on \texttt{\SthreeModelA}, a fixed projection with \texttt{expand} lowered the median LLM
API cost per run by \SthreeEOneTrimCosti{} and \SthreeEOneTrimCostiv{} on two agents with every implemented gate
passing, and when a changed task needed fields the projection had removed, the agent recovered them and they were
kept from the next run in \SthreeHFiveKept{} of \SthreeHFiveRecoveries{} cases. Offline, the projection learned with
these defaults removed \BaseLearnedStaticBigMin{} to \BaseLearnedStaticBigMax{} of what a minifier and a static
metadata filter leave on three of the five agents. The earlier studies, run with versions 0.1 and 0.2, are where
those defaults come from.
 \textbf{Study~1}, with a protocol frozen before any run, measures how much a projection
learned from other runs removes on a held-out run, where it must be learned to work, whether it transfers between
tasks, and whether the agent's answer survives when the conversation is replayed with projected tool outputs. It
finds large savings and three precise ways in which the projection fails: values the agent \emph{transforms}
before writing them (a Unix timestamp written as a date), fields the agent uses only \emph{to reason}, and, most
seriously, an agent that \emph{invents identifiers} when the identifier field has been removed. \textbf{Study~2}
adds two safeguards motivated by those failures and tests, again with a protocol frozen in advance, whether they
remove those specific losses, on the conversations that showed them, while keeping most of the saving.
\textbf{Study~3}, added in version 2 of this paper, measures in live runs, under two further frozen protocols, what
version 0.4 of the library added: a \emph{fixed} projection with \texttt{expand} (E1), learning that \emph{adapts}
while the agent works, with \texttt{expand} requests as evidence (E2), and, without trimming, a text that teaches the
agent how to organise its tool calls, on a second model (E3, reported in an appendix). It also reports offline
analyses, among them a comparison with a minifier and with a static filter fixed before any number, and post-hoc
audits of the answers, each labelled as such.

The contributions are: (i) a method to learn tool-output projections from traces, with a leave-one-out evaluation
that estimates savings and misses on runs the projection has never seen; (ii) evidence that, in these agents,
learning without the final answer misses far more of what is used (with the aggressive projection of version 0.1,
\MissesPooledShareRealCalls{} against \MissesPooledShareRealFull{} of the held-out structured calls it trims on the
three agents with real APIs), so a proxy placed only between agent and tools (an MCP gateway) learns poorly by
itself;
(iii) a characterisation of the failure modes, including identifier fabrication; (iv) safeguards and their test;
(v) an open-source implementation with a bilingual telemetry report that is useful on its own, independently of
any trimming; (vi) a live measure of that implementation's fixed trimming and adaptive learning, and of a call-organising
text, with the hypotheses that failed reported as failed; (vii) an offline comparison with minification and a
pre-registered static filter, which separates what learning adds from what any serialiser or metadata filter
removes.
